
\section{Lower Bounds in the Standard Computational Model}
\label{sec:SAT}
Our hardness result in \Cref{sec:Pi} shows that {\sc oracle Exact Matroid Basis (EMB)} is hard, leading to the unconditional lower bounds for all non-trivial oracle MOL problems in~\Cref{sec:tight_absolute}. Nonetheless, these hardness results consider matroids with general membership oracles, and do not give a lower bound for matroids that can be efficiently encoded. This is particularly important, as in some settings oracle models differ from non-oracle models w.r.t complexity~\cite{canetti2004random,chang1994random}. Moreover, some matroids show up in problems that can be encoded efficiently.
 This includes partition matroids, graphic matroids, linear matroids, etc (see, e.g., \cite{Sc03} for a survey on various families of matroids). Next, we formally define an efficient encoding of matroids.
	
	
	\begin{definition}
		\label{def:encode}
		
		A function $f:\{0,1\}^* \rightarrow 2^{\mathbb{N}} \times 2^{2^{\mathbb{N}}}$ is called {\em matroid decoder} if for every $I \in \{0,1\}^*$ it holds that $f(I) = \left(E_{f(I)}, \cI_{f(I)} \right)$ is a matroid, and the following holds. \begin{enumerate}
			\item There is an algorithm that given $I \in \{0,1\}^*$ returns $E_{f(I)}$ in time $|I|^{O(1)}$.\label{cond:1}
			\item There is an algorithm that given $I \in \{0,1\}^*$ and $S \subseteq E_{f(I)}$ decides if $S \in \cI_{f(I)}$ in time $|I|^{O(1)}$. \label{cond:2}
		\end{enumerate} 
	\end{definition} 
	
	
	
	There is a simple matroid decoder that can decode every matroid $(E,\cI)$ (such that $E\subseteq \mathbb{N}$), in which the encoding $I$ explicitly lists $\cI$. 
	However, using such a matroid decoder, the encoding size of a matroid might be very large, up to $|I|=\Omega\left(2^{|E|}\right)$, while we often seek algorithms with running times polynomial in $|E|$. 
	One way to overcome this difficulty is via the oracle model considered  in previous sections. However, our results in this model may suggest that the hardness of EMB and MOL problems is 	due to the intrinsic hardness of the oracle model.
	 Yet, there are families of matroids with very efficient encoding. For example, a {\em uniform matroid} $(E,\cI)$, where $\cI = \{S \subseteq E~|~|S| \leq k\}$, can be efficiently encoded using $I = (E,k) \in \{0,1\}^*$. Clearly, the time to decide membership of a given subset $S \subseteq E$ depends only on $|S|$ and $k$. Another example is given in \Cref{fig:M1}. 
	 
	


We start with a definition of an encoded variant of EMB. We technically define a different problem for every decoder $f$. The definition of the problem is analogous to the versions of EMB considered earlier in the paper, besides that the matroid is given via an arbitrary bit-string $I \in \{0,1\}^*$, that a matroid decoder $f$ decodes into a matroid $f(I)$. 

\problem{{\sc $f$-decoded Exact Matroid Basis} ($f$-decoded EMB)}{
	{\bf Decoder} & $f:\{0,1\}^* \rightarrow 2^\mathbb{N} \times 2^{2^{\mathbb{N}}}$ is a matroid decoder.\\
	{\bf Instance} &  $(I,c,T)$, where $I \in \{0,1\}^*$, $c:E_{f(I)} \rightarrow \mathbb{N}$, $T \in \mathbb{N}$.\\
	{\bf Solution} & A basis $S$ of the matroid $f(I)$ such that $c(S)=T$.\\
	{\bf Objective} & Decide if there is a solution.
}

As a simple example, consider the $f_{\textsf{u}}$-decoded EMB problem, for a specific matroid decoder $f_{\textsf{u}}$ that decodes uniform matroids. The matroid decoder $f_{\textsf{u}}$ interprets every $I \in \{0,1\}^*$ as $I = (E,k)$ where $E$ is a set (of numbers) and $k \in \mathbb{N}$, and returns the uniform matroid $f_{\textsf{u}}(I) = (E,\cI)$ such that $\cI = \{S \subseteq E~|~|S| \leq k\}$; clearly, $f_{\textsf{u}}$ is a matroid decoder. Thus, an instance of $f_{\textsf{u}}$-decoded EMB is a tuple $U = ((E,k),c,T)$ and a solution of $U$ is $S \subseteq E$ such that $|S| = k$ and $c(S) = T$; the goal, as before, is to decide if there is a solution. This problem is commonly known as the $k$-{\sc subset sum}.  


	\begin{figure}
	\centering
	
	\begin{tikzpicture}[thick]
		\tikzstyle{edge} = [->, line width=1pt]
		%\tikzset{edge/.style = {->,> = latex'}}
		\node (v1) at (0,2) {};
		\node (v5) at (9,5) {};
		\node (v6) at (9,4.5) {};
		\node (v7) at (9,4) {};
		
		
		\node (v2) at (1.5,3) {};
		\node at (2,1) {$\bf \textcolor{black}{U}$};
		\node (v3) at (4,2.5) {};
		\node (v4) at (7,2) {};
		\node (v5) at (7,3) {};
		\node at (7,1) {$\bf \textcolor{black}{V}$};
		\node at (12,1) {$\bf \textcolor{black}{W}$};
		\node (u0) at (11,2.5) {};
		\node (u1) at (12,2.5) {};
		\node (u2) at (13,2.5) {};
	
		
		\begin{scope}[fill opacity=0.9]
			\filldraw[fill=yellow!70] ($(v1)+(-0.5,0)$) 
			to[out=90,in=180] ($(v2) + (0,0.5)$) 
			to[out=0,in=90] ($(v3) + (1,0)$)
			to[out=270,in=0] ($(v2) + (1,-0.8)$)
			to[out=180,in=270] ($(v1)+(-0.5,0)$);
			
			
			
			\filldraw[fill=green!50] ($(v4)+(-1.5,0.2)$)
			to[out=90,in=180] ($(v4)+(1,2)$)
			to[out=0,in=90] ($(v4)+(0.6,0.3)$)
			to[out=270,in=0] ($(v4)+(1,-0.6)$)
			to[out=180,in=270] ($(v4)+(-1.5,0.2)$);
			
			
			%	\begin{scope}
				%	\filldraw[fill opacity=0.9]
				\filldraw[fill=red!70] ($(u1)+(-3,0)$) 
				to[out=90,in=180] ($(u2) + (0,0.5)$) 
				to[out=0,in=90] ($(u0) + (4,0)$)
				to[out=270,in=0] ($(u2) + (0,-0.8)$)
				to[out=180,in=270] ($(u1)+(-3,0)$);
			\end{scope}
			
			
			\filldraw[fill=brown!70] (v1) circle (0.1) node [right] {$u_1$};
			\filldraw[fill=darkgray!70] (v2) circle (0.1) node [below left] {$u_2$};
			\filldraw[fill=blue!70] (v3) circle (0.1) node [left] {$u_3$};
			\filldraw[fill=red!70] (v4) circle (0.1) node [below] {$v_1$};
			\filldraw[fill=purple!70] (v5) circle (0.1) node [below] {$v_2$};
			\filldraw[fill=white!70] (u1) circle (0.1) node [right] {$w_3$};
			\filldraw[fill=gray!70] (u2) circle (0.1) node [right] {$w_2$};
			\filldraw[fill=pink!70] (u0) circle (0.1) node [left] {$w_1$};
		\end{tikzpicture}
		\caption{\label{fig:M1} An example of a {\em partition matroid} $(E,\cI)$, which  can be efficiently encoded. The ground set is $E = \{u_1,u_2,u_3,v_1,v_2,w_1,w_2,w_3\}$, partitioned into three sets: $U,V,W$. The independent sets are all subsets of $E$ containing at most one element from $U,V$, and $W$; that is, $\cI = \{S \subseteq E~|~\forall X \in \{U,V,W\}: |S \cap X| \leq 1\}$. A simple efficient encoding of $(E,\cI)$ is $I = (E,U,V,W)$. %Clearly, $|I|$ is polynomial in $E$; in addition, 
			Membership can be decided efficiently given $I$, by checking the feasibility
			of a given set $S$ w.r.t. $U,V$ and $W$.} 
	\end{figure}

Recall that \Cref{thm:emb} asserts that oracle-EMB does not admit a pseudo-polynomial time algorithm. However, %Fully PTAS,
 this does not rule out that hypothetically, for every matroid decoder $f$ there is a pseudo-polynomial time algorithm  %Fully PTAS 
 for $f$-decoded EMB. The next result excludes this option. %Our objective is to rule out this option. 

\begin{theorem}
	\label{thm:embA}
		Assuming $\textnormal{P} \neq \textnormal{NP}$, there is a matroid decoder $f$ such that there is no algorithm  for $f$-\textnormal{decoded EMB} that for any  $f$-\textnormal{decoded EMB} instance $U =(I,c,T)$, where $n = \left|E_{f(I)}\right|+1$ and $m = c\left(E_{f(I)}\right) +1$, runs in time $ (n \cdot(T+2) \cdot m)^{O(1)}$. 
\end{theorem}

	
	
The proof of \Cref{thm:embA} is given towards the end of this section. Analogously to our hardness result for oracle MOL problems, we use \Cref{thm:embA} to give a hardness result for an encoded version of MOL problems. For every matroid decoder $f$ and every $P \in \cP$, we define a variant of the $P$-{\sc matroid optimization with a linear constraint ($P$-MOL)} problem %BM, EMI, CMB and KCM 
	in which the matroid is given via an arbitrary bit-string which a matroid decoder $f$ decodes into a matroid. 
	Formally, let $P \in \cP$, where $P = (\opt,\cF,\triangleleft)$, be the parameters of the $P$-MOL problem. For a matroid decoder $f$, we define the $f$-decoded $P$-MOL problem as follows.  
	
		\problem{ {\sc $f$-decoded $P$-matroid optimization with a linear constraint} (f-decoded $P$-MOL)}{
		{\bf Decoder} & $f:\{0,1\}^* \rightarrow 2^\mathbb{N} \times 2^{2^{\mathbb{N}}}$ is a matroid decoder.\\
		{\bf Instance} &   $(I,v,w,L)$, where $I \in \{0,1\}^*$,  $v:E_{f(I)} \rightarrow \mathbb{R}_{\geq 0}$, $w:E_{f(I)} \rightarrow\mathbb{R}_{\geq 0}$, $L \in \mathbb{R}_{\geq 0}$. \\
	%	{\bf Matroid} & $\cm_{f(I)} = \left(E_{f(I)},\cI_{ f(I)}\right)$.\\ 
		{\bf Objective} & $\opt~ v(S) \text{ s.t. } S \in \cF\left(f(I)\right),~  w(S) \triangleleft L$.
	}
	
	For example, consider the encoded version of the %{\sc budgeted matroid independent set (BM)} (or, 
	$P$-MOL for $P = (\max,\IS,\leq)$ with the matroid decoder $f_{\textsf{u}}$ that decodes uniform matroids. %In this variant, the matroid is decoded by a matroid decoder $f_{\textsf{u}}$ from a bit  $I = (E,k) \in \{0,1\}^*$. 
	An instance of the $f_{\textsf{u}}$-decoded $P$-MOL problem is a tuple $U = (I,v,w,L)$ where $I = (E,k)$ is a bit-string used for extracting the uniform matroid $f_{\textsf{u}}(I) = (E,\cI)$ such that $\cI = \{S \subseteq E~|~|S| \leq k\}$, $v$ is the value function, $w$ is the weight function, and $L$ is the bound. %By the definition of $f_{\textsf{u}}$, a 
	A {\em solution} of $U$ is $S \subseteq E$ such that $|S| \leq k$ and $w(S) \leq L$; the goal is to find a solution $S$ of maximum value $v(S)$. This problem is widely known as {\sc knapsack with cardinality constraint}.  
	
	%We will define similar variants of EMI, CMB and KCM later.
	%Hypothetically, while 
	%As for EMB, even though \Cref{thm:main} rules out a Fully PTAS for every non-trivial oracle MOL problem, this hardness result does not necessarily hold for the encoded variants of MOL problems. 
	Recall that $\cQ = \cP \setminus \{(\min,\IS,\leq)\}$ is the set of parameters for non-trivial MOL problems. Using the hardness of $f$-decoded, for some matroid decoder $f$ (details on $f$ are given towards the end of the section), we show the hardness of the $f$-decoded  variant of all non-trivial MOL problems.
	
	

	 \begin{theorem}
	 	\label{thm:mainA}
	 	Assuming $\textnormal{P} \neq \textnormal{NP}$, for any $P \in \cQ$ there is a matroid decoder $f$ such that there is no \textnormal{Fully PTAS} for $f$-\textnormal{decoded} $P$-\textnormal{MOL}. 
	 \end{theorem}
	 
	 In the remainder of this section, we prove \Cref{thm:embA} and  \Cref{thm:mainA}. The matroid decoder used in our proofs decodes a subclass of the $\Pi$-matroid family (see \Cref{sec:Pi}), in which the secret family $\Pi$ consists of the solutions for a  {\sc boolean satisfiability problem (SAT)} instance. 
	
	
	In a SAT instance $A = (V,\bar{V},\cC)$ with $n \in \mathbb{N}$ variables (in a slightly simplified notation), we are given a set $V = \{v_1, \ldots,v_n\}$ of variables, their negations $\bar{V} = \{\bar{v}_1,\ldots, \bar{v}_n\}$, and a set $\cC \subseteq 2^{V \cup \bar{V}}$ of clauses. 
	The goal is to decide if there is a set $S \subseteq [n]$ %an assignment $\sigma:V \rightarrow \{0,1\}$ 
	satisfying that for all $C \in \cC$ there is $i \in [n]$ such that one of the following holds.   \begin{itemize}
		\item $v_i \in C$ and $i \in S$.
		\item $\bar{v}_i \in C$ and $i \notin S$. 
	\end{itemize} Such a set $S$ is called a {\em solution} of $A$; let $\mathcal{S}(A)$ be the set of solutions of a SAT instance $A$. In addition, let $n(A) = n$ be the number of variables in the instance $A$. 
	The family of SAT-matroids is the subfamily of $\Pi$-matroids where $\Pi = \mathcal{S}(A)$ for some SAT instance $A$ (for the notation and definition of $\Pi$-matroids, see \Cref{def:Matroid}). Specifically, 
	%\footnote{Recall that $\Pi$-matroids are defined in \Cref{sec:Pi}.} Formally, 
	
	\begin{definition}
		\label{def:SATmatroid}
		%	Let $n,k,\alpha \in \mathbb{N}_{>0}$, and 
		Let $A$ be a \textnormal{SAT} instance, $k \in \left[n(A)\right]$, and $\alpha \in \left[  {n(A)}^{2}\right]$. Define the {\bf SAT-matroid} on $A,k,\alpha$ as $M_{n(A),k,\alpha}(\mathcal{S}(A)) = \left(\left[n(A)\right],\cI_{n(A),k,\alpha}(\mathcal{S}(A))\right)$. 
	\end{definition}
	We show below that SAT-matroids can be  encoded efficiently. %\footnote{Technically, we also restrict $k$ and $\alpha$ to be small w.r.t. $n$.} 
	For every $I \in \{0,1\}^*$, we interpret $I$ as $I = (A,k,\alpha)$, where $A$ is a \textnormal{SAT} instance, $k \in \left[n(A) \right]$, and $\alpha \in \left[ {n(A)}^2 \right]$; w.l.o.g., we may assume that every $I = (A,k,\alpha) \in \{0,1\}^*$ can be interpreted in that manner; moreover, we can assume that $n(A) \leq |A|$, where $|A|$ is the encoding size of $A$. The following definition defines a matroid decoder that decodes SAT-matroids from a bit-string. %Additionally, for simplicity we use $A_I=A$, $n_I = n(A)$, $k_I = k$, and $\alpha_I = \alpha$ for $I = (A,k,\alpha)$. 
	\begin{definition}
		\label{def:SATencoder}
		Define the {\em SAT-decoder}  as the function $f_{\textnormal{SAT}}:\{0,1\}^* \rightarrow 2^\mathbb{N} \times 2^{2^{\mathbb{N}}}$ such that for all $I = (A,k,\alpha) \in \{0,1\}^*$ it holds that $f_{\textnormal{SAT}}(I) = M_{n(A),k,\alpha}(\mathcal{S}(A))$. 
	\end{definition}
	In the next result, we show that the SAT-decoder $\fs$ is indeed a  matroid decoder.
	\begin{lemma}
		\label{lem:encodable}
		%The \textnormal{SAT-decoder} 
		$\fs$ is a matroid decoder such that $|I| = |A|^{O(1)}$ for every $I=(A,k,\alpha) \in \{0,1\}^*$. 
	\end{lemma}
	Recall the sets  $\mathcal{J}_{n,k}$,  $\mathcal{K}_{n,k,\alpha}$, $\mathcal{L}_{n,k,\alpha}(\Pi)$ and $\cI_{n,k,\alpha}(\Pi)$ were defined in  \Cref{def:Matroid}.
	\begin{proof} %[Proof of 
		We first show that $\fs$ is a matroid decoder. %; otherwise, $f$ returns the trivial matroid $(\emptyset,\{\emptyset\})$. 
		%
		%Let $n,k,\alpha \in \mathbb{N}_{>0}$, and let $A = (V,\bar{V},\cC)$ be a \textnormal{SAT} instance with $n$ variables. We show that $M_{n,k,\alpha}(\mathcal{S}(A))$ is an encodable matroid. 
		%
		For all $I = (A,k,\alpha) \in \{0,1\}^*$  it holds that $\mathcal{S}(A) \subseteq 2^{\left[n(A)\right]}$; thus, by \Cref{lem:1} it holds that $f_{\textnormal{SAT}}(I)$ returns a matroid. Moreover, for all $I = (A,k,\alpha) \in \{0,1\}^*$ it holds that $k,\alpha = {n(A)}^{O(1)}$ and $n(A) \leq |A|$ by our encoding; hence, $|I| = |A|^{O(1)}$.  It remains to prove that $\fs$ satisfies \Cref{cond:1,cond:2} in \Cref{def:encode}.
		\begin{enumerate}
			\item  	\Cref{cond:1}. Given $I = (A,k,\alpha) \in \{0,1\}^*$, we can compute $E_{f_{\textnormal{SAT}}(I)} = \left[n(A)\right]$ in time $|I|^{O(1)}$ by identifying the number of variables (i.e., $n(A)$) in $A$. 
			\item 		\Cref{cond:2}. Let $I = (A,k,\alpha) \in \{0,1\}^*$ and let $n = n(A)$; we show that we can decide membership in the matroid $\fs(I)$ in time $\io$.  %and recall that  %and let $I = (A,k,\alpha)$. %If $f(I) = (\emptyset, \{\emptyset\})$ then trivially we can decide if $S \subseteq \emptyset$ belongs to $\{\emptyset\}$ in time $\io$. Otherwise, 
			% there are $k,\alpha \in \mathbb{N}_{>0}$, and a SAT instance $A$ such that 
			%	$f_{\textnormal{SAT}}(I) = M_{n_I,k_I,\alpha_I}(\mathcal{S}(A_I))$. 
			For every $S \subseteq \left[n\right]$ we can compute $|S|$ and $\textsf{sum}(S)$ in time $|I|^{O(1)}$. Thus, validating if $S \in \mathcal{J}_{n,k}$ or $S \in \mathcal{K}_{n,k,\alpha}$ can be computed in time $\io$ regardless of $\mathcal{S}(A)$. In addition, %SAT is known to be in NP; thus, 
			deciding if $S \in \mathcal{S}(A)$ can be computed in time $|A|^{O(1)} = |I|^{O(1)}$ by confirming that all clauses are satisfied. Thus, by \eqref{eq:1} we can decide if $S \in \mathcal{L}_{n,k,\alpha}(\mathcal{S}(A))$. Overall, deciding whether $S \in \mathcal{J}_{n,k}$, $S \in \mathcal{K}_{n,k,\alpha}$, or $S \in \mathcal{L}_{n,k,\alpha}(\mathcal{S}(A))$ can be computed in time $|I|^{O(1)}$. By the above and \eqref{eq:1}, we can compute in time $\io$ correctly if a set $S$ belongs to $\cI_{\fs (I)} = \cI_{n,k,\alpha}(\mathcal{S}(A))$. 
		\end{enumerate}
		Thus, $\fs$ is a matroid decoder by \Cref{def:encode}, which gives the statement of the lemma.  %, satisfying that $|I| = |A|^{O(1)}$ for every $I=(A,k,\alpha) \in \{0,1\}^*$. . 
	\end{proof} 
	
	

	
	
	
	We will use an algorithm for $f$-decoded EMB to obtain an algorithm for SAT. To this end, consider the following family of {\em structured } $\fs$-decoded EMB instances.
		
		
											\begin{definition}
			\label{def:problemE}
			An $\fs$-{\sc decoded EMB} instance $(I,c,T)$, where $I = (A,k,\alpha)$ is called {\em structured} if  $\alpha \in \left[ {n(A)}^2 \right]$, $T = \alpha$, and $c: \left[ n(A) \right] \rightarrow \mathbb{N}$ such that for all $i \in \left[ n(A) \right]$ it holds that $c(i) = i$ %\textnormal{\textsf{sum}}(S)$.  
		\end{definition}
		
		
		The next observation immediately follows from the definition of structured $\fs$-decoded EMB instances and SAT-matroids. 
		
		\begin{observation}
			\label{obs:YesNoD}
			%For every $n,k,\alpha \in \mathbb{N}_{>0}$, $\Pi \subseteq 2^{[n]}$ 
			for any structured $\fs$-{\sc decoded EMB} instance $U = (I,c,T)$, where $I = (A,k,\alpha)$, and $S \subseteq \left[ n(A)\right]$, it holds that: $S$ is a solution for $U$ if and only if $S \in \mathcal{S}(A)$, $|S| = k$, and $\textnormal{\textsf{sum}}(S) = c(S) = \alpha$. 
		\end{observation}
		
		%The following hardness result, 
		We show that given a polynomial algorithm that decides structured $\fs$-decoded EMB instances, we can decide SAT. This result easily imply \Cref{thm:embA} as shown afterwards.  
		
%		\begin{definition}
%			An instance $s$
%		\end{definition}

		\begin{lemma}
				\label{lem:SAT}
				Assuming $\textnormal{P} \neq \textnormal{NP}$, there no algorithm that decides every $\fs$-decoded \textnormal{EMB} structured instance $(I,c,T)$ in time $\io$. 
			\end{lemma}
	\begin{proof}
	
	Assume that $\textnormal{P} \neq \textnormal{NP}$ and assume towards a contradiction that there is an algorithm $\cA$ that decides every $\fs$-decoded EMB structured instance $(I,c,T)$ in time $\io$. We describe the following algorithm $\cB$ which decides SAT using $\cA$. Let $A$ be a SAT instance and let $n = n(A)$. 
	%We describe the following algorithm $\cB$ which decides $(A,k)$. 
	\begin{enumerate}
		\item 	For every $k \in \left[n\right]$ and $\alpha \in \left[n^2\right]$, let $I_{k,\alpha} = (A,k,\alpha)$; execute $\cA$ on the structured instance $U_{k,\alpha} = (I_{k,\alpha},c,\alpha)$. %such that $f_{\textnormal{SAT}}(I_{\alpha}) = M_{n,k,\alpha}(\mathcal{S}(A))$. 
		\item  If $\emptyset$ is a solution for $A$,  return ``yes" on $A$. 
		\item If there are $k \in \left[n\right]$ and $\alpha \in \left[n^2\right]$ such that $\cA$ returns ``yes" on $U_{k,\alpha}$, %is a ``yes"-instance or $(ii)$ $\emptyset$ is a solution for $A$, %decides that $\mathcal{L}_{n,k,\alpha}(\mathcal{S}(A)) \neq \emptyset$, %yields a solution $S \in \mathcal{S}(A)$ 
	 return ``yes" on $A$. %has a solution.
	 \item Otherwise, return ``no" on $A$. %does not have a solution.
	\end{enumerate}
	To show that $\cB$ decides SAT, consider the following cases. 
	\begin{itemize}
		\item $A$ is a ``yes"-instance. Then, there is a solution $S \in \mathcal{S}(A)$. If $\emptyset \in \mathcal{S}(A)$ then $\cB$ returns ``yes" on $A$. Otherwise, there are $k \in [n]$, $\alpha \in \left[n^2\right]$, and a solution $S \in \mathcal{S}(A)$ such that $|S| = k$ and $ \textnormal{\textsf{sum}}(S) = \alpha$. Since $U_{k,\alpha}$ is structured, we have $c(S) = \alpha$; therefore, $S$ is a solution for $U_{k,\alpha}$ by \Cref{obs:YesNoD}. Consequently, $\cA$ returns ``yes" on $U_{k,\alpha}$; hence, $\cB$ returns ``yes" on $A$.  
		
	
		
		\item $A$ is a ``no"-instance. Then, $\mathcal{S}(A) = \emptyset$. %In particular, $\emptyset \notin \mathcal{S}(A)$. %for every $k \in [n]$ and $\alpha \in \left[ n^2\right]$ there is no solution for instance $I_{A,k,\alpha}$. 
		Assume towards a contradiction that $\cB$ returns ``yes" on $A$. Thus, since $\emptyset \notin \mathcal{S}(A)$, there are $k \in [n]$ and $\alpha \in \left[ n^2\right]$ such that $\cA$ returns ``yes" on $U_{k,\alpha}$, implying that there is a solution $S$ for $U_{k,\alpha}$. Since $U_{k,\alpha}$ is structured, by \Cref{obs:YesNoD} it holds that $S \in \mathcal{S}(A)$. 
		This is a contradiction that $\mathcal{S}(A) = \emptyset$. Thus, $\cB$ returns ``no" on $A$. 
	\end{itemize} 
	Hence, $\cB$ decides the SAT instance $A$ correctly. We analyze the running time below.  
	
	 By our encoding, it holds that $n \leq |A|$ as $A$ contains $n$ distinct variables. In addition, observe that the identity function $c:[n] \rightarrow [n]$ can be encoded using at most $n \cdot \log (n)$ bits (this is the weight function of every structured instance $U_{k,\alpha}$ constructed in the algorithm, for $k \in [n]$ and $\alpha \in \left[n^2\right]$). Then, for all $k \in [n]$ and $\alpha \in \left[ n^2\right]$ observe that the encoding of the instance $U_{k,\alpha}$ satisfies: $$|U_{k,\alpha}| = |A|+k+\alpha+n \cdot \log (n) \leq |A|+n+n^2 +n \cdot \log (n)= O\left(|A|^2\right).$$ %The first equality holds since the function $c$ does not need to be encoded as it simply the identity function for structured instances. 
	 Moreover, the number of iterations (i.e., the number of distinct options for the parameters $k,\alpha$) is bounded by $n \cdot n^2 = n^3 = O(|A|^3)$. Thus, the running time of each iteration is bounded by $|U_{k,\alpha}|^{O(1)} = |A|^{O(1)}$ by the running time guarantee of $\cA$. Therefore, our algorithm decides SAT in time $|A|^{O(1)}$. Since SAT is known to be NP-Hard (see, e.g., \cite{garey1997computers} for more details), we reach a contradiction to the existence of $\cA$. %unless $\textnormal{P} = \textnormal{NP}$ such an algorithm $\cA$ cannot exist. 
\end{proof}

From the above result, the hardness of the more general $\fs$-decoded EMB easily follows. As an immediate corollary, \Cref{thm:embD} gives the proof of \Cref{thm:embA}. 

	\begin{restatable}{theorem}{thmembD}
	\label{thm:embD}
	Assuming $\textnormal{P} \neq \textnormal{NP}$, there is no algorithm  for $\fs$-decoded \textnormal{EMB} that runs in time $ (n \cdot(T+2) \cdot m)^{O(1)}$, for any  $\fs$-decoded \textnormal{EMB} instance $U =(I,c,T)$ where $n = \left|E_{\fs(I)}\right|+1$ and $m = c\left(E_{\fs(I)}\right) +1$. 
\end{restatable}

\begin{proof}
	Let $U = (I,c,T)$ be a structured $\fs$-decoded EMB instance and let $I = (A,k,\alpha)$. Recall that by our encoding it holds that $|I| \geq |A| \geq n(A)$. Thus, %; let $n = |E_{f(I)}|+1$ and $m = c(E_{f(I)})+1$. 
	since $U$ is structured, it follows that  $\left|E_{\fs(I)}\right| = n(A) \leq |I|$ and 
	$$c\left(E_{\fs(I)}\right) = \textsf{sum}\left(\left[n(A)\right]\right) \leq \frac{n(A) \cdot (n(A)+1)}{2} = O\left(|I|^2\right).$$
	Thus, by the above, an algorithm $\cA$ that decides every $\fs$-EMB  instance $(I',c',T')$ in time $(n' \cdot(T'+2) \cdot m')^{O(1)}$, where  $n' = |E_{\fs(I')}|+1$ and $m' = c\left(E_{\fs(I)}\right) +1$, in particular decides $U$ in time $\io$. As $U$ is an arbitrary structured instance, unless $\textnormal{P} = \textnormal{NP}$, by \Cref{lem:SAT} such an algorithm $\cA$ cannot exist.  
\end{proof}

%We can finally prove \Cref{thm:mainD}. 
%\thmMainD*

	Finally, %we use the SAT-decoder to 
	we show that %there is an $f$-BM problem which does not have 
the variants of non-trivial matroid optimization with a linear constraint (MOL) problems, in which the decoding is performed by the SAT-decoder $\fs$, do not admit Fully PTAS under the standard assumption $\textnormal{P} \neq \textnormal{NP}$.   The proof is similar to the proof of \Cref{thm:main} in \Cref{sec:tight_absolute}. \Cref{thm:mainD} directly gives the proof of \Cref{thm:mainA}.

\begin{restatable}{lemma}{thmMainD}
	\label{thm:mainD}
	Assuming $\textnormal{P} \neq \textnormal{NP}$, for any $P \in \cQ$ there is no \textnormal{Fully PTAS} for $\fs$-\textnormal{decoded} $P$-\textnormal{MOL}. 
\end{restatable}

\begin{proof}
	%Let $P \in \cQ$ be a non-trivial MOL problem. 
	Assume that $\textnormal{P} \neq \textnormal{NP}$ and assume towards a contradiction that there is a Fully PTAS $\cA$ for $\fs$-decoded $P$-MOL. We use $\cA$ to decide $\fs$-decoded EMB. Let $U = (I,c,T)$  be an $\fs$-decoded EMB instance. %let $n = |E|+1$ and let $m = c(E)+1$; observe that $n,m \neq 0$. C
	Consider the following algorithm $\cB$ that decides $U$. In the algorithm, we use the reduction described in \Cref{def:reduction} and an error parameter defined in \eqref{eq:epsI}. 
	\begin{enumerate}
		\item Construct the $\fs$-decoded $P$-MOL instance $X = (I,v_J,w_{J,P},L_{J,P})$, where \newline $J = \left(E_{\fs(I)},\cI_{\fs (I)},c,T\right)$ is an EMB instance.\label{step:1D} 
		\item 
		Execute $\cA$ with the input $X$ and $\eps_J$ (see \eqref{eq:epsI}). \label{step:2D} 
		\item If $\cA$ returns ``no" on $X$: Return ``no" on $U$. %does not have a solution. \label{step:3D} 
		\item Otherwise, let $S \leftarrow \cA(X,\eps_J)$ be the returned solution by $\cA$.\label{step:4D} 
		\item Return ``yes" on $U$ %has a solution 
		if and only if $v_J(S) = H_J \cdot k_J+T$. \label{step:5D} 
	\end{enumerate}
	
	Observe that $U,J$ share the same matroid, weight function, and target value; the only difference between them is the representation of the matroid, that is encoded in $U$ while in $J$ the representation of the matroid is not specified. Similarly, $X$ and $R_P(J)$ also share the same matroid, value function, weight function, and bound; the only difference between them is that $X$ gives a concrete encoding of the matroid.  
	
	Let $n = |E|+1$ and let $m = c(E)+1$. Note that $X$ can be trivially constructed from $U$ in time $\left(n \cdot (T+2) \cdot m\right)^{O(1)}$ by \Cref{def:reduction}. % moreover, in time $\left(n \cdot (T+2) \cdot m\right)^{O(1)}$ we can compute the instance $X$ from $J$. %because $\frac{1}{\eps} = \left(n \cdot (T+1)\cdot m\right)^{O(1)}$; 
	Thus, the running time of $\cB$ on $U$ is $\left(n \cdot (T+2) \cdot m\right)^{O(1)}$, since $\cA$ is a Fully PTAS for $\fs$-decoded $P$-MOL and by the selection of the error parameter  $\eps_J$ in \eqref{eq:epsI}. %For the correctness,
	We now show correctness. %For the 
	
	
	\begin{itemize}
		\item If $\cB$ returns ``yes" on $U$. %has a solution. 
		Then, Step 4 of the algorithm computes a solution $S$ for $X$ satisfying $v_J(S) = H_J \cdot k_J+T$. Clearly, $S$ is a solution for $R_P(J)$ as well. By \Cref{thm:H}, %since $S$ is a solution for $R_{P}(I)$ such that $v_I(S) = H_I \cdot k_I+T$ it also holds that 
		$S$ is also a solution for $J$. As $J$ and $U$ share the same matroid, weight function, and target value, we conclude that $S$ is a solution for $U$. Therefore, $U$ is a ``yes"-instance.     
		
		
		\item If $U$ is a ``yes"-instance. Then, $J$ has a solution. Therefore, by \Cref{thm:H} it holds that $R_P(J)$ has a solution, which implies that $X$ has a solution. Thus, since $\cA$ is a Fully PTAS for $\fs$-decoded $P$-MOL, $\cA$ returns a $(1+\eps_J)$-approximate solution $S$ for $X$; clearly, $S$ is also a $(1+\eps_J)$-approximate solution for $R_P(J)$. %Thus, $S$ is a $(1+\eps_J)$-approximation for $R_P(J)$. 
		Then, by \Cref{thm:H} it follows that $v_J(S) = H_J \cdot k_J+T$. Thus, $\cB$ returns ``yes" on $U$.
	\end{itemize}  Hence, $\cB$ decides $\fs$-decoded EMB in time $\left(n \cdot (T+2) \cdot m\right)^{O(1)}$, contradicting \Cref{thm:embD}.  
\end{proof}


	
	

	
